\lesson{structure}{项目文件怎样分工和衔接}
\diagram{p_files}
文件把聊天中的要求留下来。图中以 / 结尾的是文件夹，\texttt{.md} 是可直接阅读、带简单排版的文字文件。\texttt{README.md} 说明启动和检查；\texttt{AGENTS.md} 是 Codex 等智能体支持的规则格式。Claude Code 是否直接加载它取决于版本与配置，也可在 \texttt{CLAUDE.md} 中用 \texttt{@AGENTS.md} 导入。\texttt{.specify/} 保存项目原则与流程支持文件，功能需求放在当前功能目录。

\begin{tabularx}{\linewidth}{@{}L{38mm}Y@{}}
\toprule 文件 & 负责回答的问题\\\midrule
\texttt{spec.md} & 要做什么、哪些内容不做、看到什么结果才算完成；原始要求保留原文或可回查链接\\
\texttt{research.md} & 哪些问题查过、依据在哪里、比较过什么方案、为什么选择、还有哪些未知；链接实验记录\\
\texttt{plan.md} / \texttt{tasks.md} & 前者记录实现方法；后者列出任务顺序、状态和检查方法\\
\texttt{PROGRESS.md} & 跨会话记录实际进展、未完成事项和下一次继续的位置\\
检查记录 & 实际运行了什么、针对哪个版本、结果和失败证据在哪里\\
\bottomrule\end{tabularx}
\begin{prompt}{整理项目文件与读取入口}
项目：[路径]，任务：[目标]，允许更新：[文档范围]。先定位已有规则、需求、调研、方案、任务、进度与证据；只补缺失且有用的文件，沿用 Spec Kit 当前功能目录，不另建副本。路径不可读先报缺口；无写权限返回草稿。\par
规则只补项目约定、修改范围、检查命令和文档入口；调研与实验用简短索引指向原文。保留已有有效内容，不移动或删除原件。列出更新路径及维护条件；按工具要求刷新规则，核对加载来源，无法观察时标为未验证，不把创建文件当成已生效。
\end{prompt}
\textbf{索引帮助查找，原文保留依据。}会话压缩本身不会改写磁盘上的项目文件；关键信息应提前保存，并在后续会话中可访问。文件仍在不等于 AI 已读取，恢复时按索引重读当前规格、进度及相关证据。需求变更同步方案与任务，完成后更新结果。小任务可合成短记录，不必凑齐文件。
